% id: 82-01
% book: 베이직쎈 확률과 통계 · 05 확률변수와 확률분포
% section: 개념 23 확률변수와 확률분포
% passage: 다음 확률변수 $X$가 가질 수 있는 값을 모두 구하시오.
% figure: none
% answer: 
% answer_source: 
% variant_level: 0 (원본 전사)
% 이 파일은 build.mjs 가 items.json 에서 만든다. 고칠 때는 items.json 을 고친다.
\smash{\raisebox{1.5mm}{\dmkichul{베이직쎈 확률과 통계 82쪽 01번}}}
서로 다른 3개의 동전을 동시에 던질 때, 뒷면이 나오는 동전의 개수 $X$
\dmhint{동전의 앞면을 $\mathrm{H}$, 뒷면을 $\mathrm{T}$라 할 때, 표본공간 $S$는\\[3pt] $S=\{\mathrm{HHH}{,}\mkern3mu\mathopen{}\mathrm{HHT}{,}\mkern3mu\mathopen{}\blank[2.4em]{,}\mkern3mu\mathopen{}\mathrm{HTT}{,}\mkern3mu\mathopen{}\mathrm{THH}{,}\mkern3mu\mathopen{}\blank[2.4em],$\\[3pt] $\blank[2.4em]{,}\mkern3mu\mathopen{}\blank[2.4em]\}$\\[3pt] 따라서 $X$가 가질 수 있는 값은 \quad $0$, $\blank$, $\blank$, $\blank$}
